% Insert this section before the existing "Do not regress" section in the
% Tessaris Pilot / COO Mission Layer technical handover. It uses only packages
% already present in that document's preamble.

\section*{Pilot operating team: mission execution bridge and department-local teaching}

This release connects the visible COO Mission Room to AION's governed mission
runtime. A founder can now describe a business outcome in ordinary language;
the model selected and released in the Vault proposes a bounded plan, AION
validates the proposed capabilities, and eligible work is placed in the
appropriate Department Pilot workspace. The release also moves teaching by
demonstration into each Department Computer card and makes the purpose and
state of the computer viewer explicit.

This is not a Qwen-specific implementation. Qwen, a connected cloud model, or
another accepted local model may be selected in the Vault. The model is used
for reasoning and proposal generation only. AION remains the authority for
capability admission, computer access, approvals, persistence, audit evidence
and execution receipts.

\subsection*{End-to-end mission flow}

The Mission Room now follows this sequence:

\begin{enumerate}[leftmargin=*]
\item The founder enters an outcome and selects \emph{Plan and launch
mission}. The desktop sends a \texttt{MissionRequest} to
\texttt{POST /api/aion/business/pilot-team/\{workspace\}/missions/launch}.
The request includes the desired outcome, requesting identity, constraints and
deliverables.
\item \texttt{PilotOperatingTeamService.launch\_mission()} first creates a
durable mission with a unique identity, revision, timeline, hand-off depth and
parallel-assignment limits.
\item \texttt{CooMissionService} resolves the workspace's released Vault
selection. It receives only a bounded fact projection assembled from
authoritative departmental business containers. Every projection records its
retrieval state, source identifiers and retrieval time.
\item The selected model returns a structured proposal. It may nominate only
registered Department Pilot capabilities; free-form model text is not treated
as execution authority.
\item Each valid proposal item is converted into a typed department queue
item. AION rebinds it to the durable mission identity and evaluates it through
the External Tool Gateway.
\item A sealed assignment is written into the assigned Department Computer's
task directory. The corresponding capsule is updated with the mission ID,
task ID, capability, objective, activity state and last-activity time.
\item The Mission Room refreshes from persisted state. A successfully planned
mission with assignments becomes \texttt{running}; unavailable or rejected
planning remains \texttt{queued} with an explicit reason. No silent fallback
to canned advice or a different model is permitted.
\end{enumerate}

\subsection*{Department evidence boundary}

The COO planner does not receive the whole workspace. The service projects a
limited set of authoritative containers for the relevant department and
redacts secret fields before model use or persistence.

\begin{longtable}{p{.22\linewidth}p{.70\linewidth}}
\textbf{Department} & \textbf{Authoritative projection}\\ \hline
Finance & Financial model, finance inbox and department intelligence.\\
Sales & Department intelligence and business operating model.\\
Marketing & Brand foundation and department intelligence.\\
Support & Department intelligence and operational runtime summary.\\
People & Organisation authority and business structure.\\
Operations & Business operating model, operational runtime summary and boardroom snapshot.\\
Products and services & Business map and business financial model.\\
\end{longtable}

If a source has not been published, the projection records
\texttt{not\_published}; the model must not invent the missing fact. Passwords,
passcodes, one-time codes, tokens, API keys, authorisation headers and other
secret-labelled values are recursively removed by the runtime redaction
boundary.

\subsection*{Assignment and capsule state}

Each department assignment contains the department, outcome, registered
capability, task identity, delegating actor, dependency list, queue item,
gateway decision and creation time. The assignment is sealed with an integrity
hash before it is written to the department workspace. Its state is derived
deterministically:

\begin{itemize}[leftmargin=*]
\item \texttt{waiting\_approval} when the queue item requires exact approval;
\item \texttt{queued} when gateway policy admits the work;
\item \texttt{blocked} when gateway policy refuses it.
\end{itemize}

The mission record stores the redacted model selection, proposal, approval
request, assignments and append-only timeline entry. The field
\texttt{external\_side\_effect\_executed} is initially false. Creating or
planning a mission therefore cannot be mistaken for proof that an email was
sent, a payment was made, a record was changed, or another external side
effect occurred.

\subsection*{What ``Watch computer'' means}

Every Department Pilot has a persistent browser profile, scoped files and
department grants. \emph{Watch computer} opens that isolated workspace and
passes the capsule's department ID, browser-profile ID, mission ID, task ID,
activity state and current objective to the Electron browser shell.

The viewer displays one of two explicit states:

\begin{itemize}[leftmargin=*]
\item \textbf{Idle:} ``No mission assigned.'' Watching exposes the isolated
workspace but does not invent instructions or start work.
\item \textbf{Active:} the mission strip displays the persisted activity state,
objective and task identity for the assignment bound to that capsule.
\end{itemize}

If the same browser window is reused, the shell reloads whenever the selected
profile or mission-state key changes. This fixes the earlier stale-view issue
in which taking control or switching department context could leave the viewer
showing old state. \emph{Take control} transfers the capsule controller to the
founder; \emph{Return to Pilot} restores Pilot control. Neither button bypasses
the gateway or the approval rules.

\subsection*{Department-local teaching by demonstration}

The former global teaching panel and department selector have been removed.
Each of the eight Department Computer cards now owns its own skill name,
expected result and \emph{Start teaching} control. The desktop layout renders
four cards per row, reduces to two columns below 900 pixels, and uses one
column below 560 pixels.

Teaching operates as follows:

\begin{enumerate}[leftmargin=*]
\item The founder enters a skill name and expected result on the intended
department card.
\item \emph{Start teaching} opens that department's isolated browser and starts
the native interaction recorder. Only one department recording may be active
at a time.
\item The founder demonstrates one short process. Recorded interactions are
converted into semantic steps with inputs, validation, failure behaviour and
approval boundaries; raw secrets are not retained.
\item \emph{Stop and create draft} submits the observed steps, department,
outcome, source systems and \texttt{stop\_and\_report} failure policy to the
skill compiler.
\item The result remains a draft until it has passed safe validation and
review. Publication does not grant unbounded execution authority.
\end{enumerate}

Compatibility is retained for demonstrations recorded by the older interface:
both \path{requires_approval} and \path{approval_required} compile into
the same mandatory approval stop. This closes a legacy naming mismatch that
could otherwise have weakened a recorded skill's approval boundary.

\subsection*{Taught processes on the workflow canvas}

After a demonstrated skill passes safe validation and human publication, AION
exports it as a version-pinned module in the workflow canvas's \emph{Taught
Processes} group. The same node can be inserted directly from the Demonstrated
Skill Library. Draft and failed skills are not exposed as executable modules.

The node retains the skill identity and published hash, owning department,
required inputs, output contract and approval boundaries. Canvas compilation
preserves these fields, so a saved workflow cannot silently adopt a later edit
of the demonstrated process. Connector triggers, built-in application nodes,
branches, delays and taught-process nodes can therefore be composed in one
graph.

Executing a taught-process node creates an idempotent
\path{workflow_skill_run} in the owning Department Computer. The parent AION
Flow changes to \texttt{waiting}; downstream nodes remain queued until the
Department Computer returns verified evidence and the workflow is resumed.
A successful completion requires evidence. If the process reports an external
effect at a recorded approval boundary, its evidence must also include the
exact approval receipt hash. Repeated delivery of the same idempotency key
reuses the existing run rather than repeating the browser process.

This wait-and-resume contract is deliberate. Dispatching a browser job is not
reported as completion, and the desktop renderer cannot provide arbitrary
steps or a fabricated hash as execution authority.

\subsection*{Authority and failure behaviour}

The following rules are invariant:

\begin{itemize}[leftmargin=*]
\item The selected model can propose work but cannot call connectors, control
the browser, manufacture an approval hash or directly mutate external state.
\item Capabilities must exist in AION's registered department capability set
before they can become queue items.
\item External changes remain held for exact-payload approval and must pass
through the server-side External Tool Gateway.
\item Model, provider, evidence, policy or gateway failure is fail-closed. The
mission remains persisted with its reason and no unapproved side effect.
\item Demonstrated browser replay currently performs safe validation only.
Live replay remains disabled until each action can enter the governed tool
queue, approval and receipt path.
\end{itemize}

During installed-application verification, a real multi-department mission
reached the Vault-selected provider but received HTTP 429 (provider rate
limit). The mission was correctly retained as \texttt{queued} with
\texttt{planner\_status=model\_unavailable}, zero assignments and no external
action. This is the intended fail-closed result, not permission to substitute
another model.

\subsection*{Persistence and audit record}

Operating-team state is workspace-scoped beneath the Pilot runtime store.
Mission records and department assignments are individually sealed; important
transitions also enter the hash-chained audit log. A mission timeline records
whether planning dispatched work, produced no department tasks, or was
blocked. Completion packs remain the evidence-bearing proof of finished work;
a model response, a green UI state or an open browser is not a completion
receipt.

COO delegation retains two loop limits: a maximum hand-off depth of four and a
maximum of eight parallel assignments per mission. These limits prevent
unbounded agent-to-agent delegation while still permitting cross-department
coordination.

\subsection*{Primary implementation files}

\begin{itemize}[leftmargin=*]
\item \path{backend/modules/aion_business/runtime/pilot_operating_team_service.py}:
mission creation, evidence projection, governed dispatch, task persistence,
capsule state, teaching compilation and audit records.
\item \path{backend/modules/aion_business/api/pilot_operating_team_api.py}:
typed Pilot Team endpoints, including \texttt{missions/launch}.
\item \path{backend/services/aion_mission_mode/pilot_capability_adapter.py}:
registered capability mappings used by the Department Pilot router.
\item \path{desktop/mac/src/pilot_operating_team_ui.js}: Mission Room request,
department cards, card-local teaching state and mission-status rendering.
\item \path{desktop/mac/src/pilot_operating_team.css}: four-column Department
Computer layout, responsive breakpoints, activity state and embedded teaching
controls.
\item \path{desktop/mac/electron/main.js} and
\path{desktop/mac/electron/isolated-browser.html}: isolated profile selection,
mission-state refresh and the visible idle/active mission strip.
\end{itemize}

\subsection*{Verification}

The focused service and desktop contract suite contains 19 passing tests. It
covers workspace bootstrap, capsule isolation, human takeover and return,
secret-free teaching, skill validation and publication, routines and duplicate
trigger protection, Mission Room planning and dispatch, gateway binding,
mission controls, completion packs, scoped memory, worker nodes, templates and
audit-chain verification. Python compilation, desktop JavaScript syntax,
Electron main-process syntax and isolated-browser inline-script parsing also
passed. The packaged application was installed and visually checked with eight
Department Computer cards arranged as two rows of four, local teaching controls
on every card, and the explicit idle message in the computer viewer.
